% !TeX root = ../组合数学笔记.tex
% 由主文档通过 \input 载入；请编译主文档。
\setcounter{chapter}{10} % 保留原书章号。
\chapter{匹配、染色与图的结构}
\label{ch:matching-coloring}

分配资源要求互不冲突地选择边，安排时间要求给相邻对象不同颜色。这些问题需要结构条件与构造方法共同作用。特别是，“再也不能直接加入一条边”只是一种局部停止条件，并不等于已经找到最多的边。

\section{二分图、匹配与 Hall 定理}

\begin{definition}
图 $G=(L\cup R,E)$ 若 $L,R$ 不交且每条边恰有一个端点在 $L$、一个在 $R$，称为\emph{二分图（bipartite graph）}。两两不共享端点的边集称为\emph{匹配（matching）}；每个 $L$ 中顶点都在某条匹配边上时，称该匹配饱和 $L$。对 $S\subseteq L$，记 $N(S)=\bigcup_{u\in S}N(u)$。
\end{definition}

\begin{theorem}[Hall 婚配定理]
有限二分图存在饱和 $L$ 的匹配，当且仅当对每个 $S\subseteq L$ 都有 $|N(S)|\ge|S|$。
\end{theorem}
\begin{proof}
必要性来自 $S$ 中各点必须匹配到不同邻点。充分性对 $|L|$ 归纳，$|L|=0,1$ 直接成立。若存在非空真子集 $S\subset L$ 满足 $|N(S)|=|S|$，则 $S$ 与 $N(S)$ 诱导的二分图满足 Hall 条件，由归纳可饱和 $S$。对剩余左侧点集的任意子集 $T$，由原图对 $S\cup T$ 的条件得
\[
 |N(T)\setminus N(S)|=|N(S\cup T)|-|N(S)|\ge|T|.
\]
故删去 $S\cup N(S)$ 的余图也满足条件，可饱和剩余左侧，两匹配合并即可。若不存在这样的 $S$，则每个非空真子集的邻点数至少多一。任选边 $uv$，其中 $u\in L$，删去 $u,v$ 后每个非空左侧子集至多失去一个邻点，仍满足 Hall 条件；由归纳取得匹配，再加回 $uv$。
\end{proof}

例如三名学生的候选题目集分别为 $\{a,b\},\{a\},\{b,c\}$，可取匹配 $1b,2a,3c$。若三人的候选集都是 $\{a,b\}$，每人虽都有选择，但取全部三人时 Hall 条件失败，故无法全部分配。条件必须对所有子集成立。

\section{增广路与最大匹配算法}

不能再添加一条边的匹配称为\emph{极大匹配（maximal matching）}；边数在全部匹配中最大者称为\emph{最大匹配（maximum matching）}。四点路径的中间一条边形成极大匹配，但两端的两条边构成更大匹配。

相对于匹配 $M$，一条边在 $M$ 内外交替、两个端点均未匹配的路径称为\emph{增广路（augmenting path）}。将其边的选中状态翻转后，仍为匹配，且大小增加一。反之，若另有 $|M'|>|M|$，考察对称差 $M\mathbin{\triangle}M'$；其中顶点度数至多为二，各分支为交替路径或偶圈。偶圈上两类边等多，而总数不等意味着某条路径的 $M'$ 边多一，该路径便是 $M$ 的增广路。因此匹配最大当且仅当不存在增广路。

在二分图中，从所有未匹配左侧顶点出发，沿未匹配边向右、匹配边向左搜索；到达未匹配右侧顶点便找到增广路。算法从空匹配开始，反复搜索并翻转，失败时停止。每轮增加一条匹配边，至多成功 $|V|/2$ 轮；邻接表实现的一轮搜索为 $O(|V|+|E|)$，故这一直接实现总成本为 $O(|V|(|V|+|E|))$。结论来自上述增广判据，不依赖遍历顺序。

\section{顶点覆盖与 K\H{o}nig 定理}

\emph{顶点覆盖（vertex cover）}是与每条边至少相交于一个端点的顶点集。每条匹配边都须由覆盖中的不同顶点照顾，故任意匹配大小不超过任意覆盖大小。

\begin{theorem}[K\H{o}nig 定理]
有限二分图的最大匹配大小等于最小顶点覆盖大小。
\end{theorem}
\begin{proof}
取最大匹配 $M$，按上一节的交替搜索规则，从所有未匹配左侧点出发，记可达集为 $Z$。令 $C=(L\setminus Z)\cup(R\cap Z)$。若某边 $uv$ 未被覆盖，则 $u\in L\cap Z$、$v\in R\setminus Z$；若该边不在 $M$ 中，可从 $u$ 到达 $v$，矛盾；若在 $M$ 中，已匹配的 $u$ 只能由其匹配伙伴 $v$ 到达，也矛盾。故 $C$ 为覆盖。

不存在可达的未匹配右侧点，否则得到增广路；所有未匹配左侧点都在 $Z$。每条匹配边的两个端点要么都可达，要么都不可达，因而恰有一个端点在 $C$ 中，且 $C$ 不含未匹配点。于是 $|C|=|M|$，结合一般不等式得到结论。
\end{proof}

同时给出大小相同的匹配与覆盖，就是最优性证书。二分图假设不可删除：三角形最大匹配大小为 $1$，最小顶点覆盖大小为 $2$。

\section{图染色与染色多项式}

将顶点赋予 $q$ 种有名称颜色，使邻接点异色，称为\emph{正常顶点染色（proper vertex coloring）}。最少所需颜色数称为色数 $\chi(G)$。若最大度数为 $\Delta$，依次给各点选一种尚未被已染邻点使用的颜色，则 $\Delta+1$ 种颜色总够用，因为被禁用的颜色至多为 $\Delta$ 种。这给出上界及构造，但一般不保证最少用色。

记 $P_G(q)$ 为使用 $q$ 种有名称颜色的正常染色数。若 $e=uv$ 是一条边，$G-e$ 的正常染色按 $u,v$ 是否同色分类：异色者对应 $G$ 的染色，同色者对应把 $u,v$ 合并所得 $G/e$ 的染色。缩边产生的平行边可合并，所以
\[
 P_G(q)=P_{G-e}(q)-P_{G/e}(q).
\]
从无边图的 $q^{|V|}$ 出发归纳，证明这是关于 $q$ 的多项式，称为\emph{染色多项式（chromatic polynomial）}。$n\ge1$ 的树先染根，再沿树为每个孩子避开父色，得 $q(q-1)^{n-1}$；$n\ge3$ 的圈由删边--缩边归纳得 $(q-1)^n+(-1)^n(q-1)$。计数不对颜色名称置换取商。

\section{平面图与 Euler 公式入门}

能在平面中画出且边仅在共同端点相交的图称为\emph{可平面图（planar graph）}；固定这样一个嵌入后称为平面图（plane graph）。平面被嵌入分成若干面，包括无界外部面。以下使用简单闭曲线将平面分成内外两部分的拓扑事实；这是 Jordan 曲线定理的内容，此处作为已知事实，不在组合教材中证明。

对连通平面图，若 $n,m,f$ 分别为顶点、边、面数，则
\begin{equation}
 n-m+f=2.
 \label{eq:planar-euler}
\end{equation}
证明可不断删去圈上的边：连通性不变，边两侧不同的面合并，$m,f$ 各减一，故左侧不变；最终为生成树，其 $m=n-1,f=1$，得到等式。这里关于圈边分隔面的判断正是所用平面拓扑事实。

对 $n\ge3$ 的连通简单平面图，按面边界计数时桥在同一面上计两次，每个面的边界长度至少为三，因此 $3f\le2m$。代入式\eqref{eq:planar-euler}得 $m\le3n-6$。若还是二分图，则没有奇圈，各面边界长度至少为四，得 $m\le2n-4$。不连通情形可先在各分支之间添加不相交的边使其连通，保留相应性质后应用上界。$K_5$ 有 $10>9$ 条边，$K_{3,3}$ 有 $9>8$ 条边，故均不可平面；边数满足这些上界则仅是必要条件。

\section{习题}
\begin{enumerate}
  \item\label{ex:matching-coloring-1} 三名学生的候选集合为 $\{a,b\},\{a,b\},\{b,c\}$，给出饱和匹配，并验证 Hall 条件。
  \item\label{ex:matching-coloring-2} 在四点路径上先选中间边，写出一条增广路和翻转后的匹配。
  \item\label{ex:matching-coloring-3} 计算四边形的染色多项式及使用三种有名称颜色的正常染色数。
  \item\label{ex:matching-coloring-4} 证明任意非空简单可平面图有一个度数不超过 $5$ 的顶点，并由删除顶点归纳构造六色染色。
\end{enumerate}

